\section{Related Work}
\label{sec:related_work}

Our work intersects model merging, test-time adaptation, and post-training low-bit quantization for resource-constrained edge systems.

\subsection{Linear Mode Connectivity \& Model Merging}
Linear Mode Connectivity (LMC) asserts that models sharing initialization lie within the same low-loss basin, enabling weight interpolation without catastrophic accuracy barriers \cite{frankle2018lottery, entezari2021role, garipov2018loss}. Building on this, model merging combines independent single-task experts directly in weight space.
Methods like Model Soups \cite{wortsman2022model} average homologous weights fine-tuned on the same task, while Task Arithmetic \cite{ilharco2022editing} introduces task vectors (subtracting base weights from expert weights) to linearly blend different tasks.

To resolve inter-task interference, RegMean \cite{jin2023regmean} uses activation statistics to align layers, and Fisher merging \cite{matena2022merging} uses Fisher information for parameter weighting. TIES-Merging \cite{yadav2023resolving} prunes weights and performs sign consensus to prevent cancellation. SAIM \cite{anonymous2026saim} leverages sharpness-aware minimization to encourage isotropic flatness, facilitating connectivity in deep landscapes.

\subsection{Adaptive Test-Time Merging}
Dynamic, layer-wise merging coefficients yield superior multi-task performance by accounting for layer-specific importance. AdaMerging \cite{yang2024adamerging} minimizes logit entropy over unlabeled calibration streams to find layer-wise parameters at test time, while SyMerge \cite{jung2025symerge} leverages task synergy to balance performance. OrthoMerge \cite{yang2026orthomerge} uses Riemannian manifolds to preserve energy. 

However, these techniques assume full-precision (FP32/FP16) weight spaces. Under low-bit quantization, they suffer from mathematical bottlenecks and are highly prone to the \emph{Overfitting-Optimizer Paradox} on tiny calibration sets.

\subsection{Post-Training Quantization (PTQ) \& Edge AI}
Post-Training Quantization (PTQ) compresses parameters for edge sensors \cite{gholami2021survey, jacob2018quantization} without requiring joint training data or heavy computation. Uniform PTQ maps floats to low-bit integers using scale factors. While 8-bit (INT8) uniform per-tensor scaling is common \cite{jacob2018quantization}, 4-bit (INT4) requires per-channel scaling to prevent outlier weights from collapsing entire layers \cite{lin2023awq}. PTQ schemes like GPTQ \cite{frantar2022gptq}, SmoothQuant \cite{xiao2023smoothquant}, and AdaRound \cite{nagel2020adaround} adjust weights or scales to minimize rounding error.

\subsection{Quantization-Aware Model Merging}
Integrating low-bit quantization into model merging is an active research frontier. Direct pre-quantized merging (Q-then-M) fails as rounding violates LMC, while naive post-merge quantization (M-then-Q) introduces severe noise. Q-Merge uses STE \cite{bengio2013estimating} or zero-order 1+1 ES to optimize merging coefficients under the quantization operator. Concurrent works include 1bit-Merging \cite{onebitmerging2025}, TVQ \cite{tvq2025}, and E-PMQ \cite{epmq2026}. 

Our proposed \textbf{Q-PolyMerge} stands distinct: we are the first to identify and mathematically resolve the Overfitting-Optimizer Paradox under test-time adaptation by projecting parameters onto a low-degree continuous polynomial subspace of layer depth.